%% kapa: exam=1 round=30 year=2019 period=2 form=A subject=appraisal_law subject_ko=감정평가관계법규 seq=25 number=25
%% source: Q-Net 감정평가사 기출문제 (cst00309) artlSeq=5205205 file=extracted/1008/5205205/2019년 제30회 감정평가사 1차시험 문제지/2019년 제30회 감정평가사 1차시험 문제지/2019년도 제30회 감정평가사 1차 2교시 A형.hwp
%% conversion: kordoc 4.24.51 -> tools/gen1.py (auto); review=proofread-sonnet; flags=none
\begin{question}{25}
건축법령상 건축허가에 관한 설명으로 옳지 않은 것은? (단, 조례는 고려하지 않음)
\begin{choices}
  \item 50층의 공동주택을 광역시에 건축하려면 광역시장의 허가를 받아야 한다.
  \item 자연환경을 보호하기 위하여 도지사가 지정ㆍ공고한 구역에 건축하는 3층의 숙박시설에 대하여 시장ㆍ군수가 건축허가를 하려면 도지사의 승인을 받아야 한다.
  \item 건축허가를 받으면「자연공원법」에 따른 행위허가를 받은 것으로 본다.
  \item 건축허가시 실시하는 건축물 안전영향평가는 건축물이 연면적 10만 제곱미터 이상이고 21층 이상일 것을 요건으로 한다.
  \item 2층 건축물이 건축허가 대상이라도 증축하려는 부분의 바닥면적의 합계가 80제곱미터인 경우에는 증축에 대한 건축신고를 하면 건축허가를 받은 것으로 본다.
\end{choices}
\end{question}
\answer{4}
